\documentclass[10pt,a4paper]{article}
\usepackage{amsmath,amssymb,amsthm,graphicx}
\usepackage{amsmath,amssymb,amsthm,hyperref,xfrac}
\providecommand{\eg}{e.g.\ }
\providecommand{\ie}{i.e.\ }
\usepackage[scr=rsfs,cal=euler]{mathalfa}
\multlinegap0pt
\hyphenation{accu-rate applied assump-tion assump-tions decom-pos-ing defined esti-mate esti-mates exten-ded hypothe-ses para-me-ter para-me-ters remain-der required}
\newenvironment{enumeratei}
{\bgroup\def\theenumi{\roman{enumi}}\def\theenumii{\arabic{enumii}}\begin{enumerate}}
{\end{enumerate}\egroup}
\newenvironment{enumeratea}
{\bgroup\def\theenumi{\alph{enumi}}\begin{enumerate}}
{\end{enumerate}\egroup}

\newcommand{\Psfrac}[2]{(\sfrac{#1}{#2})}
\newcommand{\psfrac}[2]{\sfrac{(#1)}{#2}}
\newcommand{\spfrac}[2]{\sfrac{#1}{(#2)}}
\newcommand{\pspfrac}[2]{\sfrac{(#1)}{(#2)}}
\newcommand{\shalf}{\sfrac{1}{2}}

\newcommand{\bnorm}[1]{\bigl\|#1\bigr\|}
\newcommand{\Bnorm}[1]{\Bigl\|#1\Bigr\|}
\newcommand{\babs}[1]{\bigl|#1\bigr|}
\newcommand{\Babs}[1]{\Bigl|#1\Bigr|}

\AtBeginDocument{%
\DeclareMathOperator{\sinn}{sin}\let\sin\sinn
\DeclareMathOperator{\coos}{cos}\let\cos\coos
}

%\usepackage{enumitem}
\usepackage{physics}
%\usepackage[foot]{amsaddr}
\usepackage{bm}
\usepackage{tikz}
\usepackage{mathtools}
\usepackage{scalerel,stackengine}
\stackMath
\usetikzlibrary{arrows.meta}
\usepackage{subcaption}
%\usepackage[style=numeric]{biblatex}
%\addbibresource{main.bib}
%\setlength{\bibitemsep}{0.3\baselineskip}

\newenvironment{mycases}
{\begin{dcases*}}
{\end{dcases*}}
\theoremstyle{plain}
\newtheorem{theorem}{Theorem}[section]
\newtheorem{lemma}[theorem]{Lemma}
\newtheorem{definition}[theorem]{Definition}
\newtheorem{proposition}[theorem]{Proposition}
\newtheorem{corollary}[theorem]{Corollary}
\newtheorem*{corollary*}{Corollary}
\theoremstyle{remark}
\newtheorem{remark}[theorem]{Remark}
\counterwithout{figure}{section}
\counterwithout{equation}{section}

\newcommand{\br}[1]{\left(#1\right)}
\newcommand{\bbr}[1]{\bigl(#1 \bigr)}
\newcommand{\Bbr}[1]{\Bigl(#1 \Bigr)}
\newcommand{\bbbr}[1]{\biggl(#1 \biggr)}
\newcommand{\tpd}{\frac{1}{(2\pi)^{d}}}
\newcommand\pig[1]{\scalerel*[2.35pt]{\big#1}{%
\ensurestackMath{\addstackgap[1.2pt]{\big#1}}}}
\newcommand\pigl[1]{\mathopen{\pig{#1}}}
\newcommand\pigr[1]{\mathclose{\pig{#1}}}
\newcommand{\ak}{\abs{k}}
\newcommand{\ap}{\abs{p}}
\newcommand{\aal}{\abs{\alpha}}
\newcommand{\sbbr}[1]{\bigl[#1\bigr]}
\newcommand{\sBbr}[1]{\Bigl[#1\Bigr]}
\newcommand{\sbbbr}[1]{\biggl[#1\biggr]}
\newcommand{\sbr}[1]{\left[#1\right]}
\newcommand{\half}{\frac{1}{2}}
\newcommand{\p}{\partial}
\newcommand{\jb}[1]{\langle#1\rangle}
\newcommand{\bset}[1]{\bigl\{#1\bigr\}}
\newcommand{\Bset}[1]{\Bigl\{#1\Bigr\}}
\newcommand{\set}[1]{\left\{#1\right\}}
\newcommand{\mytextfrac}[2]{{\textstyle\frac{#1}{#2}}}
\newcommand{\qaq}{\quad\text{and}\quad}
\newcommand{\para}{\parallel}
\newcommand{\ep}{\varepsilon}
\newcommand{\R}{\mathbb{R}}
\newcommand{\C}{\mathbb{C}}
\newcommand{\al}{\alpha}
\newcommand{\intRd}{\int_{\R^d}}
\newcommand{\intRdk}{\int_{\R^d_k}}
\newcommand{\intRdp}{\int_{\R^d_p}}
\newcommand{\lam}{\lambda}
\newcommand{\ceils}[1]{\lceil#1\rceil}
\newcommand{\Qin}{Q_{\mathrm{in}}}
\newcommand{\hin}{h_{\mathrm{in}}}
\newcommand\reallywidehat[1]{%
\savestack{\tmpbox}{\stretchto{%
\scaleto{%
\scalerel*[\widthof{\ensuremath{#1}}]{\kern.pt\mathchar"0362\kern.1pt}%
{\rule{0ex}{\textheight}}%WIDTH-LIMITED CIRCUMFLEX
}{\textheight}%
}{2.4ex}}%
\stackon[-6.9pt]{#1}{\tmpbox}%
}
\title{Uniform Penrose stability from boundary exclusion including the directional low-frequency response}
\author{Marnie Smith}
\date{doi:10.5802/jep.353; CC BY4.0}
\begin{document}
\maketitle
\noindent\textit{Editorial note.} The parameter convention delta\_0 in(0,1] is retained as the exact source356 fragment. It is the positive regularity slack needed by the directional Cauchy–Schwarz estimate and uniform tail bound. The unrelated stronger nonlinear-theorem Sobolev assumptions on source356 and radial sufficient cases near212 are not imported.\par\medskip
\noindent\textit{Editorial note.} The low-frequency limiting response L\_0(mu,omega) is an explicit hypothesis. It is not the point value L(lambda,0;hbar)=0, and cannot be dropped when applying compactness.\par\medskip
\noindent\textit{Editorial note.} The proof itself gives the bound for hbar in[0,1], including the defined classical limit; the original uniform Penrose definition requests hbar in(0,1].\par\medskip
\noindent\textit{Editorial note.} Derivatives of widehat(g) in the directional trace estimate mean Fourier-argument derivatives evaluated on the indicated line, as also fixed by the chain-rule factor in1122. No derivative of the normalized direction map is being estimated.\par\medskip
\noindent\textit{Editorial note.} No assumption of monotone/radial equilibrium from the following proposition is needed: the counted hypotheses are the original explicit real/imaginary boundary conditions.\par\medskip
\noindent\textit{Editorial scope.} One uniform stability criterion with its full local Fourier-direction and response estimates. The following smallness/radial sufficient cases1192-1254 are not counted as extra tasks; full Hartree damping is also not claimed.\par\medskip
\noindent\textit{Allowed standard prerequisites (editorial).}\begin{itemize}
\item Complex argument principle and dominated convergence. Exclude zeros for each fixed parameter and establish continuity of the compactified response; pointwise zero-freeness is not assumed to give a uniform gap.
\item Fourier Plancherel and Riemann-Lebesgue lemmas. Move between moment/Sobolev control and Fourier traces, and control large spatial frequencies.
\item Sobolev line trace and polynomial-weight ordering. Ordinary Sobolev infrastructure; the directional weighted L1 estimate1037-1056 and all response bounds remain proved obligations.
\end{itemize}\medskip
for some constant $\delta_0\in(0,1]$

\begin{definition}[Uniform Penrose condition]
An interaction kernel \(w\colon \R^d \to \R\) and steady state \(\gamma_g = (2\pi\hbar)^d g(-i\hbar\nabla)\) satisfy the uniform Penrose condition if there exists a constant \({c_{\mathrm{P}}} > 0\) such that the inequality
\begin{align}\label{Penrose}
\inf_{\substack{\Re\lambda \geq 0 \\ k \in \R^d}} \abs{1 + \widetilde{\mathscr{L}}(\lambda, k)} \geq {c_{\mathrm{P}}}
\end{align}
holds uniformly for all \(\hbar \in (0,1]\), where
\[
\widetilde{\mathscr{L}}(\lambda, k) = \frac{2}{\hbar}\, \widehat{w}(k) \int_0^\infty e^{-\lam t}\sin\bbr{\tfrac{1}{2}\hbar t \ak^2} \widehat{g}\br{kt}\,dt.
\]

\end{definition}

\subsection{Quantum operators and norms} To formulate phase-mixing and scattering estimates in the quantum setting, quantum analogues of classical derivatives and weighted Sobolev norms are introduced. The standard and double Japanese brackets are defined by $\jb{x}:=(1+\abs{x}^2)^{1/2}$ and $\jb{x,y}:=(1+\abs{x}^2+\abs{y}^2)^{1/2}$, respectively. For a multi-index $\alpha=(\alpha_1,\dots,\alpha_d)\in\mathbb N^d$, write
$\abs{\alpha}=\alpha_1+\dots+\alpha_d$ and $D^{\alpha}=\p_1^{\alpha_1}\cdots\p_d^{\alpha_d}$.

The classical weighted Sobolev norm $\norm{\cdot}_{H^\sigma_M}$ is now recalled.

\begin{definition}[Weighted Sobolev norm]
Let $\sigma\geq0$ and let $M\geq0$ be an integer. For a function $g=g(v)$, the weighted Sobolev norm $H^{\sigma}_M$ is defined as
\begin{align*}
\norm{g}_{H^{\sigma}_M}:=\sum_{\aal\leq M}\norm{v^\alpha\jb{\nabla}^\sigma g}_{L^2}.
\end{align*}
For a function $f=f(x,v)$, the weighted Sobolev norm $H^{\sigma}_{M}$ is defined as
\begin{align*}
\norm{f}_{H^{\sigma}_{M}}:=\sum_{\aal\leq M}\norm{v^\alpha \jb{\nabla_x,\nabla_v}^{\sigma}f}_{L^2},
\end{align*}where the weight applies only to the velocity variable.\end{definition}

Throughout, integrals such as $\int_x\,dx$ denote integration over $\R^d$ in the $x$ variable. The following conventions are adopted for the Fourier transform of functions $g=g(v)$ and $f=f(x,v)$:\vspace{-3pt}
\begin{align*}
\widehat g(p):=\int_v e^{-ip\cdot v}g(v)\,dv,\qquad \widehat f(k,p):=\int_x\int_ve^{-ik\cdot x}e^{-ip\cdot v}f(x,v)\,dvdx.
\end{align*}

The following lemma shows that, in the space $H^\sigma_M$, moments and derivatives can be reordered freely, allowing them to be applied in whichever order is most convenient.
\begin{lemma}[Ordering of moments]\label{lemmaHMnorm}
Let $\sigma\geq0$, let $M\in\mathbb N$, and $g\in H^\sigma_M(\R^d)$. Then the following norms are equivalent:
\begin{align*}
\norm{g}_{H^\sigma_M}:=\sum_{\abs{\alpha}\leq M}\norm{\jb{\nabla_{k}}^\sigma (k^\alpha g)}_{L^2}\sim\sum_{\abs{\alpha}\leq M}\norm{k^\alpha\jb{\nabla_{k}}^\sigma g}_{L^2}.
\end{align*}
\end{lemma}

The Sobolev $L^2$ trace lemma is also recalled.
\begin{lemma}[$L^2$ trace {\cite[Lem.\,2.15]{BedrossianMasmoudiMouhot2018}}] \label{L2trace}
Let $g\in H^s(\R^d)$ with $s>\psfrac{d-1}{2}$ and $C\subset\R^d$ be an arbitrary straight line. Then, there holds
\begin{align*}
\norm{g}_{L^2(C)}\lesssim_{s,d}\norm{g}_{H^s(\R^d)}.
\end{align*}
\end{lemma}

\begin{align}
\widetilde {\mathscr L}(\lam,k;\hbar)
= \frac{2}{\hbar\ak} \widehat w(k) \int_0^\infty e^{-\lam \sfrac{s}{\ak}} \sin(\mytextfrac{\hbar}{2} \ak s) \widehat g\bbr{\mytextfrac{k}{\ak}s} \, ds. \label{Lingoodcoordinates}
\end{align}
To compare with the classical setting, consider the limit
\[
\widetilde{\mathscr{L}}(\lambda,k;0) := \lim_{\hbar \to 0} \widetilde{\mathscr{L}}(\lambda,k;\hbar),
\]
which evaluates to\vspace*{-3pt}
\[
\widetilde {\mathscr{L}}(\lam,k;0)
= \widehat w(k) \int_0^\infty e^{-\lam \sfrac{s}{\ak}} s \widehat g\bbr{\mytextfrac{k}{\ak}s} \, ds.
\]

Next, using the argument principle, a condition on the behaviour of \( \widetilde{\mathscr L} \) along the imaginary axis is shown to imply the uniform Penrose condition. The low frequencies require separate attention: \( \widetilde{\mathscr L}(\lam,0;\hbar)=0 \), whereas \eqref{Lingoodcoordinates} and dominated convergence give, for every $\hbar\in[0,1]$, $\Re\mu\geq0$ and $\omega\in\mathbb S^{d-1}$,\vspace*{-3pt}
\begin{align*}
\lim_{r\downarrow0}\widetilde{\mathscr L}(r\mu,r\omega;\hbar)=\widetilde{\mathscr L}_0(\mu,\omega):=\widehat w(0)\int_0^\infty e^{-\mu s}s\,\widehat g(\omega s)\,ds,
\end{align*}
so \( \widetilde{\mathscr L} \) is not in general continuous at $(\lam,k)=(0,0)$.

\begin{lemma}\label{argumentprinciplelemma}
Suppose \( w \in L^1(\R^d) \), \( g \in H^{3/2+\delta_0}_{1+\lceil d/2 \rceil}(\R^d) \), and that the following implications hold:\begin{align}
\label{impliesPenrose}\forall\hbar\in[0,1],\,\tau\in\R,\,k\in\R^d,\quad \Im\widetilde {\mathscr L}(i\tau,k;\hbar)=0\implies \Re\widetilde {\mathscr L}(i\tau,k;\hbar)>-1,
\end{align}and, for the low-frequency response,
\begin{align}
\label{impliesPenrose0}\forall\tau\in\R,\,\omega\in\mathbb S^{d-1},\quad \Im\widetilde {\mathscr L}_0(i\tau,\omega)=0\implies \Re\widetilde {\mathscr L}_0(i\tau,\omega)>-1.
\end{align} Then the uniform Penrose condition \eqref{Penrose} is satisfied.
\end{lemma}

These lemmas imply the following estimate for directional integrals of derivatives of \( \widehat{g} \):
\begin{lemma}\label{gnorm}
Let \( n \geq 0 \), and let \( \alpha \) be a multi-index. Then, for every $k\neq0$,
\begin{align*}
\int_0^\infty\jb{s}^{n}\babs{\p^\alpha_k\widehat g\bbr{\mytextfrac{k}{\abs{k}}s}}\,ds\lesssim_{d,
\delta_0,n,\aal}\norm{g}_{H^{n+{{1/2}}+\delta_0}_{\aal+\ceils{d/2}}}.\end{align*}
\end{lemma}
\begin{proof}
Applying Cauchy--Schwarz followed by Lemma~\ref{L2trace},
\begin{align*}
\int_0^\infty\jb{s}^{n}\babs{\p^\alpha_k\widehat g\bbr{\mytextfrac{k}{\abs{k}}s}}\,ds&\lesssim_{\delta_0}\br{\int_0^\infty\jb{s}^{2n+1+2\delta_0}\babs{\p^\alpha_k\widehat g\bbr{\mytextfrac{k}{\abs{k}}s}}^2\,ds}^{{1/2}}\\
&\lesssim_{d}\bbbr{\intRdk\abs{\jb{\nabla_k}^{\ceils{d/2}}\bbr{\jb{k}^{n+{{1/2}}+\delta_0}\p^\alpha_k\widehat g(k)}}^2\,dk}^{{1/2}}.
\end{align*}
Then, by Plancherel and Lemma~\ref{lemmaHMnorm}:
\begin{align*}
\int_0^\infty\jb{s}^{n}\babs{\p^\alpha_k\widehat g\bbr{\mytextfrac{k}{\abs{k}}s}}\,ds&\lesssim_d\bnorm{\jb{x}^{\ceils{d/2}}\jb{\nabla_x}^{n+{{1/2}}+\delta_0}\br{x^\alpha g(x)}}_{L^2}\\
&\lesssim\sum_{\abs{\beta}\leq\ceils{d/2}}\bnorm{x^\beta\jb{\nabla_x}^{n+{{1/2}}+\delta_0}\br{x^\alpha g(x)}}_{L^2}\\
&\lesssim_n\sum_{\abs{\beta}\leq\ceils{d/2}}\bnorm{\jb{\nabla_x}^{n+{{1/2}}+\delta_0}\br{x^{\beta+\alpha} g(x)}}_{L^2}\\
&\lesssim_\alpha\norm{g}_{H_{\aal+\ceils{d/2}}^{n+{{1/2}}+\delta_0}}.\qedhere
\end{align*}
\end{proof}

\begin{lemma}\label{boundedness and decay lemma}
Suppose $w\in L^1(\R^d)$ and $g\in H^{3/2+\delta_0}_{\ceils{d/2}}(\R^d)$. Then:\vspace*{-3pt}
\begin{enumeratea}
\item \label{(a)} As $\ak\to\infty$, the quantity $\widetilde {\mathscr L}(\lam,k;\hbar)\to0$, with decay uniform in $\Re\lam\geq0$ and $\hbar\in[0,1]$.

\item \label{(b)}There exists a constant $C_1=C_1(d,\delta_0)$, independent of $\hbar\in [0,1]$, such that for all $\Re\lam\geq0$ and $k\in\R^d$,\vspace*{-3pt}
\begin{align*}
\babs{\widetilde {\mathscr L}(\lam,k;\hbar)}\leq C_1\norm{w}_{L^1}\norm{g}_{H^{3/2+\delta_0}_{\ceils{d/2}}}.
\end{align*}
\item \label{(c)} If in addition $g\in H^{3/2+\delta_0}_{1+\ceils{d/2}}(\R^d)$, then there exists a constant $C_2=C_2(d,\delta_0)$, independent of $\hbar\in [0,1]$, such that for all $\Re\lam\geq0$ with $\lam\neq0$ and $k\in\R^d$,\vspace*{-3pt}
\begin{align*}
\babs{\widetilde {\mathscr L}(\lam,k;\hbar)}\leq C_2\frac{\ak}{\abs{\lam}}\norm{w}_{L^1}\norm{g}_{H^{3/2+\delta_0}_{1+\ceils{d/2}}}.
\end{align*}
\end{enumeratea}
\end{lemma}
\begin{proof}
First observe that \( \widetilde{\mathscr L}(\lam,k=0;\hbar) = 0 \) for all \( \Re \lam \geq 0 \) and \( \hbar \in [0,1] \). Now fix \( \hbar \in (0,1] \). For \( k \neq 0 \) and \( \Re \lam \geq 0 \), the representation \eqref{Lingoodcoordinates} together with Lemma~\ref{gnorm} yields the estimate\vspace*{-3pt}
\begin{align*}
\babs{\widetilde{\mathscr L}(\lam,k;\hbar)}
&\lesssim_d \abs{\widehat w(k)} \frac{1}{\hbar \ak} \int_0^\infty \babs{\sin\br{\mytextfrac{\hbar}{2} \ak s} \, \widehat g\bbr{\mytextfrac{k}{\ak}s}}\,ds \\
&\lesssim \abs{\widehat w(k)} \int_0^\infty s\, \babs{\widehat g\bbr{\mytextfrac{k}{\ak}s}}\,ds \\
&\lesssim_{d,\delta_0} \abs{\widehat w(k)} \norm{g}_{H^{3/2+\delta_0}_{\ceils{d/2}}}.
\end{align*}
Since \( w \in L^1(\R^d) \), the Fourier transform \( \widehat w(k) \to 0 \) as \( \abs{k} \to \infty \), establishing part~\eqref{(a)}. Moreover, \( \abs{\widehat w(k)} \leq \norm{w}_{L^1} \) for all \( k \), yielding part~\eqref{(b)}.

For part~\eqref{(c)}, the expression \eqref{Lingoodcoordinates} can be integrated by parts in \( s \), with no boundary contributions, to obtain\vspace*{-1pt}
\begin{align*}
\widetilde{\mathscr L}&(\lam,k;\hbar)
\\[-3pt]
&= 2 \widehat w(k) \frac{1}{\hbar \ak} \Bigl(-\frac{\ak}{\lam}\Bigr) \int_0^\infty \p_s \bigl( e^{-\lam \sfrac{s}{\ak}} \bigr) \sin(\mytextfrac{\hbar}{2} \ak s) \, \widehat g\bbr{\mytextfrac{k}{\ak}s} \, ds \\
&= 2 \widehat w(k) \frac{1}{\hbar \lam} \int_0^\infty e^{-\lam \sfrac{s}{\ak}} \left[ \mytextfrac{\hbar}{2} \ak \cos(\mytextfrac{\hbar}{2} \ak s) \widehat g\bbr{\mytextfrac{k}{\ak}s} + \sin(\mytextfrac{\hbar}{2} \ak s) \mytextfrac{k}{\ak} \cdot \nabla_k \widehat g\bbr{\mytextfrac{k}{\ak}s} \right] ds.
\end{align*}
Therefore,\vspace*{-5pt}
\begin{align*}
\abs{\widetilde{\mathscr L}(\lam,k;\hbar)}
&\lesssim \frac{\ak}{\abs{\lam}} \norm{w}_{L^1} \int_0^\infty \Bigl[ \babs{\widehat g\bbr{\mytextfrac{k}{\ak}s}} + s \babs{\nabla_k \widehat g\bbr{\mytextfrac{k}{\ak}s}} \Bigr] ds \\
&\lesssim_{d,\delta_0} \frac{\ak}{\abs{\lam}} \norm{w}_{L^1} \norm{g}_{H^{3/2+\delta_0}_{1+\ceils{d/2}}}.
\end{align*}
The same bounds hold for the classical limit \( \widetilde{\mathscr L}(\lam,k;0) \), with identical estimates.
\end{proof}

\begin{proof}
For each $\hbar\in[0,1]$ and $k\in\R^d$, the assumptions on $w$ and $g$ ensure that the map $\lam\mapsto \widetilde {\mathscr L}(\lam,k;\hbar)$ is analytic in the half-plane $\set{\Re\lam>0}$, continuous on $\set{\Re\lam\geq0}$ and, by Lemma~\ref{boundedness and decay lemma}~\eqref{(c)}, tends to zero as $\abs{\lam}\to\infty$. By the argument principle, the number of zeros of the function $\lam\mapsto 1+\widetilde {\mathscr L}(\lam,k;\hbar)$ in $\set{\Re\lam>0}$ is equal to the winding number of the curve\vspace*{-3pt}
\begin{align*}
\Gamma:=\bigl\{\widetilde {\mathscr L}(i\tau,k;\hbar):\tau\in \R\bigr\}
\end{align*}
around the point $-1$. The implication \eqref{impliesPenrose} ensures that $\Gamma$ does not meet the half-line $(-\infty,-1]$. It follows that $1+\widetilde {\mathscr L}(\lam,k;\hbar)$ has no zeros in $\set{\Re\lam\geq0}$. The same argument, applied to $\mu\mapsto1+\widetilde{\mathscr L}_0(\mu,\omega)$ for fixed $\omega\in\mathbb S^{d-1}$ and using \eqref{impliesPenrose0}, shows that $1+\widetilde{\mathscr L}_0(\mu,\omega)$ has no zeros in $\set{\Re\mu\geq0}$.

It remains to show that the infimum in \eqref{Penrose} is positive. For $k\neq0$, $\hbar\in[0,1]$ and $\Re\lam\geq0$, the representation \eqref{Lingoodcoordinates} reads\vspace*{-3pt}
\begin{align*}
\widetilde{\mathscr L}(\lam,k;\hbar)&\hphantom{:}=\widehat w(k)\,\Phi\br{\sfrac{\lam}{\ak},\hbar\ak,\sfrac{k}{\ak}},\\
\Phi(\mu,h,\omega)&:=\int_0^\infty e^{-\mu s}\,\frac{2}{h}\sin\br{\sfrac{hs}{2}}\widehat g(\omega s)\,ds,
\end{align*}
where the factor $2h^{-1}\sin(hs/2)$ is understood as $s$ when $h=0$, so that $\widetilde{\mathscr L}_0(\mu,\omega)=\widehat w(0)\Phi(\mu,0,\omega)$. The function $\Phi$ is continuous on $\set{\Re\mu\geq0}\times[0,\infty)\times\mathbb S^{d-1}$. Indeed, the integrand is continuous and bounded by $s\norm{\widehat g}_{L^\infty}$, so the integral over $[0,R]$ is continuous for every $R>0$, while $\abs{2h^{-1}\sin(hs/2)}\leq s$ and the Cauchy--Schwarz and trace estimates in the proof of Lemma~\ref{gnorm} give, uniformly in $\Re\mu\geq0$, $h\geq0$ and $\omega\in\mathbb S^{d-1}$,\vspace*{-3pt}
\begin{align*}
\int_R^\infty\Babs{e^{-\mu s}\,\frac{2}{h}\sin\br{\sfrac{hs}{2}}\widehat g(\omega s)}\,ds
&\leq\br{\int_R^\infty\hspace*{-1mm}\frac{s^2}{\jb{s}^{3+2\delta_0}}\,ds}^{\sfrac{1}{2}}\br{\int_0^\infty\hspace*{-1mm}\jb{s}^{3+2\delta_0}\abs{\widehat g(\omega s)}^2\,ds}^{\sfrac{1}{2}}\\
&\lesssim_{d,\delta_0} R^{-\delta_0}\norm{g}_{H^{3/2+\delta_0}_{\ceils{d/2}}},
\end{align*}
so $\Phi$ is a uniform limit of continuous functions. Now set\vspace*{-3pt}
\[
F(\mu,r,\omega,\hbar):=1+\widehat w(r\omega)\,\Phi(\mu,\hbar r,\omega),
\]
so that $F(\mu,r,\omega,\hbar)=1+\widetilde{\mathscr L}(r\mu,r\omega;\hbar)$ for $r>0$ and $F(\mu,0,\omega,\hbar)=1+\widetilde{\mathscr L}_0(\mu,\omega)$.

By Lemma~\ref{boundedness and decay lemma}, there exist constants \( K > 0 \) and \( M > 0 \) such that, for \( k\neq0 \) and \( \Re \lam \geq 0 \), if either (i) \( \ak \geq K \), or (ii) \( \abs{\lam} \geq M\ak \), then
\[
\babs{1 + \widetilde{\mathscr L}(\lam,k;\hbar)} \geq \frac{1}{2}
\]uniformly in $\hbar\in[0,1]$, while $\babs{1 + \widetilde{\mathscr L}(\lam,0;\hbar)}=1$. Writing $\lam=r\mu$ and $k=r\omega$, the remaining region corresponds to the compact set \[
\mathcal K:=\left\{
(\mu, r, \omega, \hbar) : \,
\Re\mu\geq0,\,\abs{\mu} \leq M,\,
0\leq r \leq K,\,
\omega\in\mathbb S^{d-1},\,
\hbar \in [0,1]
\right\}.
\]
On this compact set, the function $F$ is continuous, since $\widehat w$ and $\Phi$ are, and it does not vanish, by the first paragraph. Hence
\[\textstyle
{c_{\mathrm{P}}} := \min\set{\sfrac{1}{2},\,\min_{\mathcal K}\abs{F}} > 0
\]
is a lower bound for $\babs{1 + \widetilde{\mathscr L}(\lam,k;\hbar)}$ over $\Re \lam \geq 0$, $k\in\R^d$ and $\hbar \in [0,1]$.
This establishes the uniform Penrose condition \eqref{Penrose}.
\end{proof}

\begin{thebibliography}{99}
\bibitem{BedrossianMasmoudiMouhot2018} J. Bedrossian, N. Masmoudi \& C. Mouhot - “Landau damping in finite regularity for unconfined systems with screened interactions”, Comm. Pure Appl. Math. 71 (2018) no. 3, p. 537-576
\end{thebibliography}
\end{document}
